\section{The arithmetic interpretation}
\label{sec:basis}

\subsection{What a presentation means}

Use the repository's ordinary intuitionistic primitive recursive arithmetic.
The only object sort is the natural numbers. Each fixed primitive recursive
(PR) definition has its own symbol. The proof calculus admits first-order
intuitionistic deductions and induction on quantifier-free formulas. In the Lean
implementation the symbols are \texttt{PR k}, terms are \texttt{PTm n}, and the
judgment with assumptions $\Gamma$ and conclusion $A$ is \texttt{Prf PRAAx}
$\Gamma\ A$. Here $n$ counts free numerical variables. We give schemata in fixed
PR algorithms, allowing their numerical parameters, including finite tables, to
vary. A schema can be instantiated by substituting terms and the corresponding
correctness derivations. Its quantification over algorithms belongs to the
metatheory.

A rational can be a signed numerator and a positive denominator. For a total
natural coding, represent the integer numerator by a difference of naturals and
the denominator by one plus a natural. Equality and order are cross-products;
an unreduced fraction is a valid result. Encode finite lists by length and
contents, or use the explicit coding in Section~\ref{sec:oracle}. Each
operation used below has a PR loop bound. One may use the same finite integer
kernel for real, vector, and table observations. Fixed finite product distances
can use the maximum of coordinate distances.

In a rational pseudometric basis $(D,d)$ let $e_p=2^{-p}$. A name is a family
$x_p$ satisfying $d(x_p,x_q)\le e_p+e_q$. Equality of two names means
$d(x_p,y_q)\le e_p+e_q$ at all $p,q$. This definition only needs basis distances,
not a prior construction of an infinite sequence as a set inside arithmetic.
For fixed PR families these are first-order formulas. Externally the same
conditions describe arbitrary Cauchy names in the completion. A bound
$|x-r|\le b$ for a named scalar and a rational means
$|x_p-r|\le b+e_p$ for every $p$. Sums, rational multiples and order inequalities
of named scalars are read with the sum of the individual absolute error bounds.
This convention also applies to the $F(q)$ and $L(f)$ notation below.

An elementary rule removes auxiliary precision parameters. Suppose rational
$A,B$ satisfy $A\le B+c2^{-r}$ for every $r$, where $c\ge0$ is a supplied
rational. If $A>B$, the positive difference has an explicitly positive
numerator. With $A-B=a/d$ and $c=b/s$, choose, for instance, $r=bd+1$;
then $2^r>bd$ and $c2^{-r}<a/d$. Decidable rational order gives $A\le B$.
All products and powers are PR. This allows calculations with named scalars to
be expanded into rational inequalities one observed precision at a time.

\subsection{A rule for transferring estimates}

\begin{theorem}[Two-run error rule]\label{thm:tworuns}
Suppose the stage-$p$ algorithm returns a basis element $v$ and a finite
transcript. There is a PR producer on the admitted numerical inputs. Assume
that from any two valid transcripts at stages $p,q$ on compatible inputs one
can derive
\[
                          d(v,w)\le e_p+e_q.
\]
Then the stage algorithm names an output point, and equivalent inputs have
equivalent outputs. If two proposed regular names $y,z$ satisfy
$d(y_i,v_k)\le e_i+\rho(k)$ and $d(z_j,v_k)\le e_j+\sigma(k)$ for all
$i,j,k$, with a supplied schedule $k(t)$ such that
$\rho(k(t))+\sigma(k(t))\le e_t$, then $y$ and $z$ are equal as names.
\end{theorem}
\begin{proof}
Use the two-run inequality on the same input, at two precisions, to get
regularity. Apply it to two names of the same input to get extensionality.
For uniqueness fix one precision of each of the two proposed outputs. Compare
both to an arbitrarily fine observation of the produced point. Add the
two bounds by the triangle inequality, then eliminate the surplus term with
the preceding rational rule. Each use instantiates a universal numerical
hypothesis. There is no selection of infinitely many witnesses.
\end{proof}

One common way of proving the two-run hypothesis is to compare each run to
the same rational test functional with error $e_p$. In the law construction
the distinguishing test is itself produced from the two finite outputs.
In the root construction one compares rational brackets before estimating
a residual. A checker which accepts candidate observations supplies a soundness
implication. To obtain the theorem's producer one must also give the candidate,
an explicit bounded search with a proved success premise, or a PR iteration
with a proved terminal condition.

A finite proof that uses a loop should state its invariant at a fixed final
precision. For a comparison needed at every $i<N$, a bounded conjunction is
compiled as the product of Boolean comparisons, or as the test that the sum of
failure bits is zero. The induction formula is an equation of terms. A universal
contract may be held as a hypothesis and instantiated at a computed index
inside a step. Inducting directly on a universal statement about all functions,
all levels, or all future inputs would require a different rule.

\subsection{Relation to the ProveIt packets}

The development uses revision \texttt{b2341af5370c} of ProveIt. Its
finite-query normal form already codes a bounded finite path, and its continuous
frameworks already give completion, effective compactness and residual-to-name
rules. Our use of finite observations builds on those results. Three more
specific comparisons determine the applications.

\begin{itemize}
\item The analysis library has polynomial bisection, root names and uniqueness
results. Section~\ref{sec:conditioned-zero} asks what comparison at a midpoint can
be replaced by when only a noisy observation of a general function is available.
\item The calculus packet proves FTC II assuming derivative information.
\texttt{TaggedFTC} proves the summation of upper and lower cell estimates.
Section~\ref{sec:dyadicapplication} constructs derivative information from a
condition on the value tables, including behavior at old grid knots.
\item The expected-value packet defines a simple-law name using finite
couplings as data. Section~\ref{sec:laws} produces such couplings from
marginal tables on the line and, one step earlier, produces the marginal
tables from an oracle for expectations of continuous polygonal probes.
\end{itemize}
Revision-specific links and reproduction instructions appear in the final
section and in the source inventory in the archive.
